\begin{center}
	\section*{Homework 10 - 3.10, 3.11}
	Due May 1 \\
	Uzair Hamed Mohammed
\end{center}

% 4, 10, 12
% 4, 8, 12, 16, 20

\subsection*{3.10 Order Statistics}
\begin{enumerate}
	\item[4.] A random sample of size 5 is drawn from the pdf \(f_Y(y) = 2y, 0 \leq y \leq 1\). Calculate \(P(Y_1^{\prime} < 0.6 < Y_5^{\prime})\).

	      \underline{Sol}: \\

	      \[
		      F_Y(y)=y^2,\quad 0\le y\le1.
	      \]
	      \[
		      P(Y_1'<0.6<Y_5') = 1 - [F(0.6)]^5 - [1-F(0.6)]^5 = 1 - (0.36)^5 - (0.64)^5.
	      \]
	      \[
		      0.36=\frac{9}{25},\quad 0.64=\frac{16}{25}.
	      \]
	      \[
		      1-\left(\frac{9}{25}\right)^5-\left(\frac{16}{25}\right)^5 = 1-\frac{59049+1048576}{9765625}=1-\frac{1107625}{9765625}=\frac{8658000}{9765625}=\frac{69264}{78125}.
	      \]
	      \[
		      \boxed{\frac{69264}{78125}}
	      \]

	\item[10.] Suppose that n observations are chosen at random from a continuous pdf \(f_Y(y)\). What is the probability that the last observation recorded will be the smallest number in the entire sample?

	      \underline{Sol}: \\

	      \[
		      \boxed{\frac{1}{n}}
	      \]

	\item[12.] Consider a system containing n components, where the lifetimes of the components are independent random variables and each has pdf \(f_Y (y) = \lambda e^{- \lambda y}, y > 0\). Show that the average time elapsing before the first component failure occurs is \(1 / n \lambda\).

	      \underline{Sol}: \\

	      Let \(Y_1,\dots,Y_n \stackrel{\text{i.i.d.}}{\sim} \text{Exp}(\lambda)\).
	      The time until the first failure is \(Y_{(1)} = \min(Y_1,\dots,Y_n)\).
	      For i.i.d. exponentials, \(Y_{(1)} \sim \text{Exp}(n\lambda)\).
	      Hence \(\mathbb{E}[Y_{(1)}] = \frac{1}{n\lambda}\). \(\square\)
\end{enumerate}

\subsection*{3.11 Conditional Densities}
\begin{enumerate}
	\item[4.] Five cards are dealt from a standard poker deck. Let X be the number of aces received, and Y the number of kinds. Compute \(P(X = 2 | Y = 2)\).

	      \underline{Sol}: \\

	      \[
		      P(X=2\mid Y=2)=\frac{\binom{4}{2}\binom{12}{1}\binom{4}{3}}{\binom{13}{2}\bigl[\binom{4}{1}\binom{4}{4}+\binom{2}{1}\binom{4}{3}\binom{4}{2}\bigr]}=\boxed{\frac{6}{91}}
	      \]

	\item[8.] In Question 3.11.7, define the random variable W to be the "majority" of x, y, and z. For example, \(W(2, 2, 1) = 2\) and \(W (1, 1, 1) = 1\). Find the pdf of \(W|_x\). For reference, Question 3.11.7 goes:

	      Suppose X, Y, and Z have a trivariate distribution described by the joint pdf
	      \[
		      p_{X, Y, Z} (x, y, z) = \frac{xy + xz + yz}{54}
	      \]

	      where x, y, and z can be 1 or 2. Tabulate the joint conditional pdf of X and Y given each of the two values of z.

	      \underline{Sol}: \\

	      Givem joint pmf \(p(x,y,z) = \frac{xy+xz+yz}{54}\) for \(x,y,z\in\{1,2\}\).

	      Compute all probabilities:

	      \[
		      \begin{array}{c|c}
			      (x,y,z) & p     \\ \hline
			      (1,1,1) & 3/54  \\
			      (1,1,2) & 5/54  \\
			      (1,2,1) & 5/54  \\
			      (1,2,2) & 8/54  \\
			      (2,1,1) & 5/54  \\
			      (2,1,2) & 8/54  \\
			      (2,2,1) & 8/54  \\
			      (2,2,2) & 12/54
		      \end{array}
	      \]

	      Marginals: \(P(X=1)=\frac{21}{54}=\frac{7}{18}\), \(P(X=2)=\frac{33}{54}=\frac{11}{18}\).

	      For \(X=1\):
	      - \(W=1\) when \((1,1,1)\), \((1,1,2)\), \((1,2,1)\): total probability \(\frac{3+5+5}{54}=\frac{13}{54}\).
	      - \(W=2\) when \((1,2,2)\): \(\frac{8}{54}\).
	      Thus \(P(W=1\mid X=1)=\frac{13/54}{21/54}=\frac{13}{21}\), \(P(W=2\mid X=1)=\frac{8}{21}\).

	      For \(X=2\):
	      - \(W=1\) when \((2,1,1)\): \(\frac{5}{54}\).
	      - \(W=2\) when \((2,1,2)\), \((2,2,1)\), \((2,2,2)\): \(\frac{8+8+12}{54}=\frac{28}{54}\).
	      Thus \(P(W=1\mid X=2)=\frac{5/54}{33/54}=\frac{5}{33}\), \(P(W=2\mid X=2)=\frac{28}{33}\).

	      \[
		      \boxed{P(W=1\mid X=1)=\frac{13}{21},\; P(W=2\mid X=1)=\frac{8}{21},\; P(W=1\mid X=2)=\frac{5}{33},\; P(W=2\mid X=2)=\frac{28}{33}}
	      \]

	\item[12.] Given the joint pdf
	      \[
		      f_{X, Y} (x, y) = 2e^{-(x+y)}, 0 \leq x \leq y, y \geq 0
	      \]

	      find
	      \begin{enumerate}
		      \item \(P(Y < 1 | X < 1)\)
		      \item \(P(Y < 1 | X = 1)\)
		      \item \(f_{Y|_x} (y)\)
		      \item \(E(Y |_x)\)
	      \end{enumerate}

	      \underline{Sol}: \\

	      Marginal of \(X\): \(f_X(x)=\int_{y=x}^\infty 2e^{-(x+y)}dy=2e^{-x}\int_x^\infty e^{-y}dy=2e^{-2x},\;x\ge0\).

	      (a) \(P(Y<1\mid X<1)=\dfrac{P(Y<1,X<1)}{P(X<1)}\).

	      \(P(X<1)=\int_0^1 2e^{-2x}dx=1-e^{-2}\).

	      \(P(Y<1,X<1)=\int_{x=0}^1\int_{y=x}^1 2e^{-(x+y)}dy\,dx=\int_0^1 2e^{-x}(e^{-x}-e^{-1})dx\)
	      \(=\int_0^1(2e^{-2x}-2e^{-x-1})dx=[-e^{-2x}+2e^{-x-1}]_0^1\)
	      \(=(-e^{-2}+2e^{-2})-(-1+2e^{-1})=e^{-2}+1-2e^{-1}=(1-e^{-1})^2\).

	      Thus \(P(Y<1\mid X<1)=\dfrac{(1-e^{-1})^2}{1-e^{-2}}=\dfrac{1-e^{-1}}{1+e^{-1}}=\dfrac{e-1}{e+1}\).

	      (b) \(P(Y<1\mid X=1)=0\) because \(Y\ge X\) almost surely, and \(X=1\) forces \(Y\ge1\).

	      (c) \(f_{Y\mid X}(y\mid x)=\dfrac{f_{X,Y}(x,y)}{f_X(x)}=\dfrac{2e^{-(x+y)}}{2e^{-2x}}=e^{-(y-x)},\quad y\ge x\).

	      (d) \(\mathbb{E}[Y\mid X=x]=\int_x^\infty y\,e^{-(y-x)}dy\). Substitute \(u=y-x\):
	      \(=\int_0^\infty (x+u)e^{-u}du=x\int_0^\infty e^{-u}du+\int_0^\infty u e^{-u}du=x\cdot1+1=x+1\).

	      \[
		      \boxed{\text{(a) }\frac{e-1}{e+1}\quad\text{(b) }0\quad\text{(c) }e^{-(y-x)},\,y\ge x\quad\text{(d) }x+1}
	      \]

	\item[16.] Suppose that X and Y are distributed according to the joint pdf
	      \[
		      f_{X, Y}(x, y) = \frac{2}{5} \cdot (2x + 3y), 0 \leq x \leq 1, 0 \leq y \leq 1
	      \]

	      find
	      \begin{enumerate}
		      \item \(f_X (x)\)
		      \item \(f_{Y|_x} (y)\)
		      \item \(P(\frac{1}{4} \leq Y \leq \frac{3}{4} | X = \frac{1}{2})\)
		      \item \(E(Y |_x)\)
	      \end{enumerate}

	      \underline{Sol}: \\

	      Given joint pdf \(f_{X,Y}(x,y) = \frac{2}{5}(2x+3y)\) for \(0 \le x \le 1\), \(0 \le y \le 1\).

	      (a) Marginal pdf of \(X\):
	      \[
		      f_X(x) = \int_0^1 \frac{2}{5}(2x+3y)\,dy = \frac{2}{5}\left[2xy + \frac{3}{2}y^2\right]_{y=0}^{1} = \frac{2}{5}\left(2x + \frac{3}{2}\right) = \frac{4x+3}{5},\quad 0\le x\le1.
	      \]

	      (b) Conditional pdf of \(Y\) given \(X=x\):
	      \[
		      f_{Y|X}(y|x) = \frac{f_{X,Y}(x,y)}{f_X(x)} = \frac{\frac{2}{5}(2x+3y)}{\frac{4x+3}{5}} = \frac{2(2x+3y)}{4x+3} = \frac{4x+6y}{4x+3},\quad 0\le y\le1.
	      \]

	      (c) For \(x=\frac12\):
	      \[
		      f_{Y|X}\!\left(y\,\middle|\,\frac12\right) = \frac{4\cdot\frac12+6y}{4\cdot\frac12+3} = \frac{2+6y}{5} = \frac{2}{5} + \frac{6}{5}y.
	      \]
	      Then
	      \[
		      P\!\left(\frac14 \le Y \le \frac34 \mid X=\frac12\right) = \int_{1/4}^{3/4} \left(\frac{2}{5}+\frac{6}{5}y\right)dy = \frac{1}{5}\left[2y+3y^2\right]_{1/4}^{3/4}.
	      \]
	      At \(y=3/4\): \(2\cdot\frac34+3\cdot\frac{9}{16} = \frac32+\frac{27}{16} = \frac{24+27}{16}=\frac{51}{16}\).
	      At \(y=1/4\): \(2\cdot\frac14+3\cdot\frac{1}{16} = \frac12+\frac{3}{16} = \frac{8+3}{16}=\frac{11}{16}\).
	      Difference: \(\frac{51-11}{16}=\frac{40}{16}=\frac52\). Then \(\frac15\cdot\frac52 = \frac12\).

	      (d) Conditional expectation:
	      \[
		      \mathbb{E}[Y\mid X=x] = \int_0^1 y\cdot\frac{4x+6y}{4x+3}\,dy = \frac{1}{4x+3}\int_0^1 (4xy+6y^2)\,dy = \frac{1}{4x+3}\left[2xy^2+2y^3\right]_{0}^{1} = \frac{2x+2}{4x+3}= \frac{2(x+1)}{4x+3}.
	      \]

	      \[
		      \boxed{f_X(x)=\frac{4x+3}{5}}\quad
		      \boxed{f_{Y|X}(y|x)=\frac{4x+6y}{4x+3}}\quad
		      \boxed{P\!\left(\frac14\le Y\le\frac34\mid X=\frac12\right)=\frac12}\quad
		      \boxed{\mathbb{E}[Y\mid X=x]=\frac{2(x+1)}{4x+3}}
	      \]

	\item[20.] Suppose the random variables X and Y are jointly distributed according to the pdf
	      \[
		      f_{X, Y} (x, y) = \frac{6}{7}(x^2 + \frac{xy}{2}), 0 \leq x \leq 1, 0 \leq y \leq 2)
	      \]

	      find
	      \begin{enumerate}
		      \item \(f_X (x)\)
		      \item \(P(X > 2Y)\)
		      \item \(P(Y > 1 | X > \frac{1}{2})\)
	      \end{enumerate}

	      \underline{Sol}: \\

	      Given joint pdf \(f_{X,Y}(x,y) = \frac{6}{7}\left(x^2 + \frac{xy}{2}\right)\) for \(0 \le x \le 1,\; 0 \le y \le 2\).

	      (a) Marginal \(f_X(x)\):
	      \[
		      f_X(x) = \int_0^2 \frac{6}{7}\left(x^2 + \frac{xy}{2}\right) dy = \frac{6}{7}\left[ x^2 y + \frac{x}{2}\cdot\frac{y^2}{2} \right]_0^2 = \frac{6}{7}\left(2x^2 + \frac{x}{2}\cdot\frac{4}{2}\right) = \frac{6}{7}(2x^2 + x).
	      \]

	      (b) \(P(X > 2Y) = \int_{x=0}^1 \int_{y=0}^{x/2} \frac{6}{7}\left(x^2 + \frac{xy}{2}\right) dy\,dx\).
	      Inner integral:
	      \[
		      \int_0^{x/2} \left(x^2 + \frac{xy}{2}\right) dy = x^2\cdot\frac{x}{2} + \frac{x}{2}\cdot\frac{(x/2)^2}{2} = \frac{x^3}{2} + \frac{x}{2}\cdot\frac{x^2}{8} = \frac{x^3}{2} + \frac{x^3}{16} = \frac{9x^3}{16}.
	      \]
	      Multiply by \(\frac{6}{7}\): \(\frac{6}{7}\cdot\frac{9}{16}x^3 = \frac{54}{112}x^3 = \frac{27}{56}x^3\). Then
	      \[
		      P = \frac{27}{56}\int_0^1 x^3 dx = \frac{27}{56}\cdot\frac{1}{4} = \frac{27}{224}.
	      \]

	      (c) \(P(Y>1 \mid X>\frac12) = \frac{P(Y>1,\,X>\frac12)}{P(X>\frac12)}\).

	      First, \(P(X>\frac12) = \int_{1/2}^1 f_X(x)dx = \int_{1/2}^1 \frac{6}{7}(2x^2+x)dx = \frac{6}{7}\left[ \frac{2}{3}x^3 + \frac{1}{2}x^2 \right]_{1/2}^1\).
	      At \(x=1\): \(\frac{2}{3}+\frac12 = \frac{7}{6}\). At \(x=\frac12\): \(\frac{2}{3}\cdot\frac18 + \frac12\cdot\frac14 = \frac{1}{12}+\frac18 = \frac{5}{24}\). Difference: \(\frac{7}{6}-\frac{5}{24} = \frac{28}{24}-\frac{5}{24}=\frac{23}{24}\). Then \(P(X>\frac12)=\frac{6}{7}\cdot\frac{23}{24}=\frac{23}{28}\).

	      Numerator: \(P(Y>1,X>\frac12) = \int_{x=1/2}^1 \int_{y=1}^2 \frac{6}{7}\left(x^2+\frac{xy}{2}\right) dy\,dx\).
	      Inner integral:
	      \[
		      \int_1^2 \left(x^2+\frac{xy}{2}\right) dy = x^2(2-1) + \frac{x}{2}\cdot\frac{2^2-1^2}{2} = x^2 + \frac{x}{2}\cdot\frac{3}{2} = x^2 + \frac{3x}{4}.
	      \]
	      Then \(\frac{6}{7}\int_{1/2}^1 \left(x^2+\frac{3x}{4}\right)dx = \frac{6}{7}\left[ \frac{x^3}{3} + \frac{3}{4}\cdot\frac{x^2}{2} \right]_{1/2}^1 = \frac{6}{7}\left[ \frac{x^3}{3} + \frac{3x^2}{8} \right]_{1/2}^1\).
	      At \(x=1\): \(\frac13+\frac38 = \frac{8}{24}+\frac{9}{24}=\frac{17}{24}\). At \(x=\frac12\): \(\frac{1}{3}\cdot\frac18 + \frac{3}{8}\cdot\frac14 = \frac{1}{24}+\frac{3}{32} = \frac{4}{96}+\frac{9}{96}=\frac{13}{96}\). Difference: \(\frac{17}{24}-\frac{13}{96} = \frac{68}{96}-\frac{13}{96}=\frac{55}{96}\). Multiply: \(\frac{6}{7}\cdot\frac{55}{96} = \frac{330}{672}=\frac{55}{112}\).
	      Thus \(P(Y>1\mid X>\frac12) = \frac{55/112}{23/28} = \frac{55}{112}\cdot\frac{28}{23} = \frac{55}{4\cdot23} = \frac{55}{92}\).

	      \[
		      \boxed{f_X(x)=\frac{6}{7}(2x^2+x)}\qquad
		      \boxed{P(X>2Y)=\frac{27}{224}}\qquad
		      \boxed{P(Y>1\mid X>\frac12)=\frac{55}{92}}
	      \]
\end{enumerate}
